\subsection{训练时的内存核算}

对一个普通的前馈模型，训练时要同时考虑：

\begin{table}[H]
\centering
\begin{tabular}{lcc}
\toprule
\textbf{部分} & \textbf{数量级} & \textbf{说明} \\
\midrule
参数 & \(N_p\) & 模型权重 \\
激活值 & \(N_a\) & 前向过程中保存的中间结果 \\
梯度 & \(N_p\) & 每个参数对应一个梯度 \\
优化器状态 & \(N_p\) 或更多 & 取决于算法 \\
\bottomrule
\end{tabular}
\caption{训练步的常见内存组成}
\end{table}

若用 float32 朴素训练，则总内存可以粗略写成：
\[
4 \times (N_p + N_a + N_p + N_p) = 4(3N_p + N_a).
\]
这也是为什么“模型参数看起来装得下”，但“训练时还是爆显存”。

